% Compile with: pdflatex sedenion_prior_work
%               bibtex sedenion_prior_work
%               pdflatex sedenion_prior_work
%               pdflatex sedenion_prior_work
% Keep sedenion_prior_work.bib in the same directory.
% Alternatively: latexmk -pdf sedenion_prior_work.tex
% The review body can be reused in another manuscript with the cited .bib entries.
\documentclass[11pt,a4paper]{article}
\usepackage[T1]{fontenc}
\usepackage[utf8]{inputenc}
\usepackage{newtxtext}
\usepackage{amsmath,amssymb,mathtools}
\usepackage{newtxmath}
\usepackage{microtype}
\usepackage[a4paper,top=25mm,bottom=26mm,left=27mm,right=27mm]{geometry}
\usepackage[numbers,sort&compress]{natbib}
\usepackage{xurl}
\usepackage[hidelinks]{hyperref}
\usepackage{bookmark}
\hypersetup{
  pdftitle={Associators, Norm Defects and Zero Divisors: Prior Work},
  pdfsubject={A focused literature review of Cayley--Dickson algebra and related physical proposals},
  pdfkeywords={octonions, sedenions, associator, zero divisors, norm defect, Hopf construction, literature review}
}
\newcommand{\R}{\mathbb R}
\newcommand{\C}{\mathbb C}
\newcommand{\OO}{\mathbb O}
\newcommand{\SSS}{\mathbb S}
\newcommand{\ip}[2]{\langle #1,#2\rangle}
\newcommand{\DN}{\Delta_N}
\DeclareMathOperator{\Rea}{Re}
\DeclareMathOperator{\Ann}{Ann}
\DeclareMathOperator{\spec}{spec}
\numberwithin{equation}{section}
\setlength{\parskip}{3pt plus 1pt}
\setlength{\parindent}{1.2em}
\setlength{\emergencystretch}{2em}
\setlength{\bibsep}{6pt plus 1pt}
\allowdisplaybreaks[1]
\raggedbottom
\pagestyle{plain}
\title{\textbf{Associators, Norm Defects\\and Zero Divisors}\\[5pt]
  \large Prior work on Cayley--Dickson algebras\\and related physical proposals}
\author{}
\date{25 September 2026}

\begin{document}
\maketitle
\vspace{-1.5em}
\begin{abstract}
This focused review summarises prior work relevant to the identity linking
an octonionic associator to the failure of norm composition in the real
sedenions. It covers Moreno's associator operators and higher Hopf
construction, the annihilator and spectral results of Biss and
collaborators, and Morita's equivalent scalar cross-term calculation.
Gillard--Gresnigt's generation construction and Tang's later associator-mass
proposal are treated separately as physical applications or hypotheses.
The central algebraic relationship has clear antecedents; a self-contained
rederivation should not be presented as a newly discovered identity.
No priority claim for a first occurrence, or validation of a physical mass
interpretation, is made here.
\end{abstract}

\section{Scope and conventions}

The review concerns the eight works in the accompanying bibliography.
Journal articles and the specific preprint versions consulted are
identified separately. Numbered-result references below refer to the
arXiv versions specified in the bibliography, except for Morita's
publisher version and Tang's version~1. The notation ``norm defect'' is
used here to compare the sources; it is not assumed to be their shared
terminology.

To fix signs, write $x=(a,b)$ and $y=(c,d)$ with $a,b,c,d\in\OO$, and use
\begin{equation}
 (a,b)(c,d)=(ac-\bar d\,b,\;da+b\bar c),
 \qquad N(x)=|a|^2+|b|^2.
 \label{eq:product}
\end{equation}
Here $N$ is the squared positive Euclidean norm on the \emph{real}
sedenions. Define
\begin{equation}
 \DN(x,y)=N(x)N(y)-N(xy),\qquad
 [a,b,c]=(ab)c-a(bc),\qquad
 \ip{u}{v}=\Rea(u\bar v).
 \label{eq:definitions}
\end{equation}
The identity under discussion is
\begin{equation}
 \boxed{\DN((a,b),(c,d))=-2\ip{[a,b,c]}{d}.}
 \label{eq:central}
\end{equation}
Section~\ref{sec:morita} explains its equivalence to Morita's published
cross-term calculation. The convention changes and elementary
rearrangements displayed in this review are explanatory comparisons,
not quotations of formulas in identical notation from every source.

\section{Moreno: associators and multiplication operators}

\subsection{Early zero-divisor work: 1997 preprint, 1998 article}

\citet{Moreno1998ZeroDivisors} studies zero divisors in real
Cayley--Dickson algebras through associators and left-multiplication
operators. The introduction explicitly connects failure of norm
multiplicativity, loss of alternativity, and the appearance of zero
divisors beyond the octonions. The sedenion case receives a detailed
geometric treatment; its zero-divisor characterisation is subsequently
cited as Corollary~2.14 by Biss and collaborators.
This supplies an early direct precedent for the relationship, rather
than merely background on nonassociative algebras.
\citep{Moreno1998ZeroDivisors,BissDuggerIsaksen2008Annihilators}

\subsection{An explicit associator block: 2005 preprint}

In \emph{Constructing zero divisors}, Moreno's Theorem~2.3 writes the
square of multiplication by a doubled element in block form. For
nonzero pure elements $a,b$ in the lower algebra, set
\begin{equation}
 \mathcal A=L_a^2+R_b^2,\qquad
 S(z)=[a,z,b]=R_bL_a(z)-L_aR_b(z).
\end{equation}
Then
\begin{equation}
 L_{(a,b)}^2=
 \begin{pmatrix}\mathcal A&-S\\S&\mathcal A\end{pmatrix}.
 \label{eq:moreno}
\end{equation}
Thus the associator is an explicit off-diagonal operator in the doubled
algebra. Moreno uses such operators and their spectra to construct and
analyse zero divisors.~\citep[Sec.~2, Theorem~2.3; Sec.~3]{Moreno2005Constructing}

\paragraph{Comparison with the scalar defect.}
When $a,b\in\OO$ are pure, octonionic alternativity gives
$\mathcal A=-(|a|^2+|b|^2)I$.
Equation~\eqref{eq:moreno} therefore isolates the associator contribution
relative to the norm-preserving diagonal term. This is the operator
counterpart of evaluating the scalar defect for a particular second
input; it does not assign that defect a physical meaning.

\section{Morita: the equivalent scalar cross-term calculation}
\label{sec:morita}

Morita's \emph{Quasi-associativity and Cayley--Dickson algebras} gives the
closest explicit scalar calculation in this review.
His Section~5 uses
\begin{equation}
 \begin{aligned}
 U&=X+e_{15}Y, &V&=Z+e_{15}W,\\
 UV&=XZ-W\bar Y+e_{15}(\bar XW+ZY),
 \end{aligned}
\end{equation}
where $X,Y,Z,W$ are octonions. In the unnumbered calculation after
Eq.~(53), the squared norm appearing inside the Study-like determinant
is
\begin{equation}
 \begin{aligned}
 |UV|^2&=(|X|^2+|Y|^2)(|Z|^2+|W|^2)
          -2\Rea(r'-s'),\\
 r'&=(XZ)(Y\bar W),\qquad
 s'=(ZY)(\bar WX).
 \end{aligned}
 \label{eq:morita}
\end{equation}
Morita attributes the surviving correction to octonionic
nonassociativity and explicitly relates failure of norm composition to
zero divisors.~\citep[Sec.~5, Eq.~(52a) and the calculation after Eq.~(53)]{Morita2014QuasiAssociativity}

\paragraph{Translation into the present convention.}
Set
\begin{equation}
 X=a,\qquad Y=\bar b,\qquad Z=c,\qquad W=\bar d.
\end{equation}
Equation~\eqref{eq:morita} becomes
\begin{align}
 \DN(x,y)
 &=2\Rea\bigl((ac)(\bar b d)-(c\bar b)(da)\bigr)\notag\\
 &=2\bigl(\ip{ac}{\bar d b}-\ip{da}{b\bar c}\bigr)\notag\\
 &=-2\ip{[d,a,c]}{b}
  =-2\ip{[a,b,c]}{d}.
 \label{eq:translation}
\end{align}
The last equality uses the alternating four-form
$\ip{[a,b,c]}{d}$.
This demonstrates equivalence, not merely thematic similarity.
Presenting the compact associator form is a useful reformulation, but
not a basis for claiming a new norm-composition identity.

\section{Biss--Dugger--Isaksen: annihilators and zero-divisor geometry}

\citet{BissDuggerIsaksen2008Annihilators} study
\begin{equation}
 \Ann(x)=\{y:xy=0\}
\end{equation}
and relate these subspaces to associators and anti-associators.
Sections~8 and~11 develop those tools; Section~12 specialises to the
sedenions. For a nonzero $(a,b)\in\SSS$, Proposition~12.1 gives
\begin{equation}
 (a,b)\text{ is a zero divisor}
 \quad\Longleftrightarrow\quad
 \Rea(a)=\Rea(b)=0,\quad
 \ip{a}{b}=0,\quad |a|=|b|.
 \label{eq:zero-condition}
\end{equation}
A nontrivial sedenion annihilator has real dimension four.
The unit zero-divisor locus is homeomorphic to the Stiefel manifold
$V_2(\R^7)$, whereas the space of ordered unit pairs with zero product
is homeomorphic to $G_2$. These are different spaces and should not be
conflated.~\citep[Sec.~12, Proposition~12.1 and following discussion]{BissDuggerIsaksen2008Annihilators}

This moves the discussion beyond isolated multiplication-table examples:
it identifies the geometry of both admissible zero divisors and their
partners. The classification is algebraic; no generation or mass
interpretation is part of this result.

\section{Biss--Christensen--Dugger--Isaksen: norm distortion as eigentheory}

\citet{BissEtAl2009Eigentheory} introduce, for $x\neq0$,
\begin{equation}
 M_x=\frac{1}{N(x)}L_{\bar x}L_x
     =\frac{1}{N(x)}L_x^{\mathsf T}L_x.
 \label{eq:M}
\end{equation}
This is a positive-semidefinite, self-adjoint real operator. Their
``eigenvalues of $x$'' mean eigenvalues of $M_x$, not of $L_x$ itself.
\citep[Definition~3.1 and Sec.~3]{BissEtAl2009Eigentheory}

\paragraph{Comparison with the defect.}
Taking the inner product of Eq.~\eqref{eq:M} with $y$ directly gives
\begin{equation}
 \DN(x,y)=N(x)\ip{y}{(I-M_x)y}.
 \label{eq:defect-operator}
\end{equation}
For a nonzero eigenvector $M_xy=\lambda y$, consequently,
\begin{equation}
 \frac{\DN(x,y)}{N(x)N(y)}=1-\lambda.
\end{equation}
Thus $\lambda=0$ detects zero product, $0<\lambda<1$ norm loss,
$\lambda=1$ norm preservation, and $\lambda>1$ norm excess.

For imaginary, linearly independent octonions $a,b$, Corollary~7.3 gives
\begin{equation}
 \spec M_{(a,b)}=\{1,1+\eta,1-\eta\},\qquad
 \eta=\frac{2|a||b|\sin\theta}{|a|^2+|b|^2},
 \label{eq:spectrum}
\end{equation}
where $\theta$ is the angle between $a$ and $b$. The multiplicities are
$8,4,4$. Equal-length orthogonal inputs yield
$\{0,1,2\}$, with multiplicities $4,8,4$, respectively
(Proposition~7.4).~\citep[Corollary~7.3 and Proposition~7.4]{BissEtAl2009Eigentheory}
This provides an existing spectral framework for the signed defect, not
a derived physical mass spectrum.

\section{Moreno: the higher Hopf construction and its zero set}

Moreno's 2004 preprint studies
\begin{equation}
 F_n(x,y)=(2xy,\;|y|^2-|x|^2),
 \qquad x,y\in A_n,
 \label{eq:hopf-moreno}
\end{equation}
where $A_n$ is the real $2^n$-dimensional Cayley--Dickson algebra.
For the composition-algebra cases, norm multiplicativity makes the
restriction to a unit sphere sphere-valued. For $n>3$, the paper instead
investigates the zero set of this construction.~\citep[Abstract and Introduction]{Moreno2004Hopf}

\paragraph{Comparison with the present Hopf-map formula.}
For sedenions, the alternative convention
\begin{equation}
 H(x,y)=(N(x)-N(y),\;2x\bar y)
\end{equation}
is obtained from $F_4$ by conjugating the second input and changing output
coordinate order and sign. Direct expansion gives
\begin{equation}
 |H(x,y)|^2=(N(x)+N(y))^2-4\DN(x,\bar y).
\end{equation}
Equation~\eqref{eq:central} implies
$\DN(x,\bar y)=\DN(x,y)$, so normalized input satisfies
\begin{equation}
 N(x)+N(y)=1
 \quad\Longrightarrow\quad
 |H(x,y)|^2=1-4\DN(x,y).
\end{equation}
This rewriting measures the obstruction by the signed defect.
The higher construction and its zero-divisor obstruction themselves
have a direct antecedent in Moreno's paper; these equations do not
produce a new sphere fibration.

\section{Physical applications and speculative mass proposals}

\subsection{Gillard--Gresnigt: three fermion generations}

\citet{GillardGresnigt2019Generations} work with the
\emph{complexified} sedenions $\C\otimes\SSS$.
After selecting a special direction and constructing a primitive
idempotent, the authors identify three complex-octonion subalgebras
sharing a quaternionic subalgebra. Their associated Clifford-algebra
ideals describe three fermion generations with unbroken
$SU(3)_c\times U(1)_{\mathrm{em}}$ symmetry.
The selection and projector are part of the construction; this should
not be recast as the claim that an arbitrary sedenion presentation
contains only three octonionic subalgebras.
\citep[Abstract]{GillardGresnigt2019Generations}

This is complementary representation-theoretic work, rather than a
source for Eq.~\eqref{eq:central} or a norm-defect mass law. A proposed
bridge to that law must also specify the norm: the positive real norm
used here and its complex-bilinear extension are not the same object
as a positive Hermitian norm on complex amplitudes.

\subsection{Tang: an explicit associator-norm mass hypothesis}

Tang's 2025 preprint explicitly proposes a relation between lepton mass
and an associator norm in Section~5.1, Eq.~(24):
\begin{equation}
 m_{\ell}\ \propto\
 \bigl\|[\overline{\Psi},A_{\mu},\Psi]\bigr\|.
\end{equation}
The author places this within a proposed nonassociative reformulation
of quantum electrodynamics and claims generation-dependent masses and
magnetic-moment corrections. The version reviewed, posted
7 November 2025, is explicitly marked as \emph{not peer-reviewed}.
\citep[Sec.~5.1, Eq.~(24)]{Tang2025AssociatorMass}

It establishes a precedent for proposing associator norms as mass
quantities, not validation of the claimed dynamics or numerical
predictions. Its displayed hypothesis is also not identical to the
signed four-linear defect in Eq.~\eqref{eq:central}. This review reports
the proposal without auditing its physical or mathematical consistency.

\section{Implications for attribution and further work}

The sources support a clear distinction. Associator-controlled
multiplication, sedenion norm nonmultiplicativity, zero-divisor geometry,
and the obstruction in the higher Hopf construction already belong to
a developed mathematical literature. The scalar, operator, and
geometric descriptions above are compatible ways of organising that
work. They should not be conflated with later hypotheses assigning
physical mass to an algebraic quantity.
\citep{Moreno2005Constructing,Morita2014QuasiAssociativity,
BissDuggerIsaksen2008Annihilators,BissEtAl2009Eigentheory,Moreno2004Hopf}

\par\penalty-100
A defensible manuscript statement is:\nopagebreak
\begin{quote}
We use established Cayley--Dickson algebra to express the failure of
sedenion norm composition as a signed scalar projection of an
octonionic associator. The calculation is equivalent to published
cross-term expressions and is related to existing operator and
annihilator descriptions. Any identification of a quantity derived
from this defect with fermion mass is an additional physical
hypothesis.
\end{quote}

A physical extension would have to specify the observable, its
transformation law, positivity or sign prescription, dimensional scale,
and dynamics. The bibliography records relevant precedents, not an
exhaustive priority search. Neither the existence of the algebraic
identity nor the citation of an earlier mass proposal establishes a
new physical prediction.

\bibliographystyle{unsrtnat}
\bibliography{sedenion_prior_work}
\end{document}
